% Code and data availability; acknowledgement of the code review.
% Source: ULXspectra-mc docs/for_paper_session.md Note 14 (e1de75a), 2026-11-24.

\section*{Acknowledgements}
\new{We thank \textbf{[names]} for reviewing the code before publication.}
% AUTHOR: names of the colleagues who review the code (private review, Note 14).

\section*{Code and data availability}
\new{The model and the fitting code, \code{ULXspectra}, are available at \url{https://github.com/gvlipunova/ULXspectra}; the version used in this paper is archived at Zenodo (doi:10.5281/zenodo.\textbf{XXXXXXX}).
The repository includes a guide to the code (\code{docs/REVIEWER\_GUIDE.md}) and regression tests.
\textbf{[Data: the fitted SEDs and the best-fit spectra are available in the repository / on request.]}}
% Note 14: the repository is private for now; reviewers are collaborators and see tag review-2026-11-24 on branch
%   mc-transfer-corrections (the paper's code; main is older), or a git bundle of that tag. docs/REVIEWER_GUIDE.md (e1de75a):
%   where the C model, interface, sampler and figure scripts are, how to build and run the regression test (13/13 pass
%   with the C28 model after the references were recomputed). Before publication: make the repository public, merge the
%   paper's branch into main or link a release, tag the paper's version, archive it on Zenodo for the DOI.
% AUTHOR: (1) keep the placeholder until the release exists; (2) data: the raw SEDs are not in the repository, but the
%   fitted 0.3-10 keV points of 29 best refits are (summaries/best_refit_spectra/*/best_fit_spectrum.txt). Decide whether
%   they are public (and whether the reviewers may see them now), then choose the wording of the data sentence.
